\begin{table*}[tp]
\centering
\renewcommand{\arraystretch}{1.3}
\caption{Funciones Asignadas al Firmware (Software)}
\label{tab:funciones_firmware}
\begin{tabularx}{\textwidth}{|X|X|}
\hline
\textbf{Función Software} & \textbf{Módulo Firmware} \\
\hline
Lectura periódica de temperatura y humedad (DHT11) & readSensors() --- timer 5 s \\
\hline
Lectura periódica de humedad del sustrato (YL-100) & readSensors() --- timer 10 s \\
\hline
Validación de lecturas (rango válido, NaN, error) & readSensors() --- validación inline \\
\hline
Lógica de control de riego con histéresis & controlIrrigation() \\
\hline
Lógica de control térmico con histéresis & controlTemperature() \\
\hline
Actualización del estado de actuadores (GPIO) & updateActuators() \\
\hline
Serialización de datos en JSON & publishIoTData() --- ArduinoJson \\
\hline
Publicación de telemetría en MQTT / Firebase & publishIoTData() --- PubSubClient \\
\hline
Conexión Wi-Fi y reconexión automática & setup() + loop() --- WiFi.begin() \\
\hline
Temporización no bloqueante (sin delay()) & millis() + lastTime\_X \\
\hline
Constantes de umbrales y pines & config.h --- HUMIDITY\_MIN, TEMP\_MAX, \ldots \\
\hline
\end{tabularx}
\end{table*}